\chapter{The data model: a Markov chain}
\label{ch:markov}

\section{Definition}

A sequence $a = (a_1,\dots,a_{\nsites})$ is a \emph{first-order Markov chain} if, given the
present, the future is independent of the past:
\begin{equation}
\label{eq:markov-property}
P(a_i \mid a_{i-1}, a_{i-2},\dots,a_1) \;=\; P(a_i \mid a_{i-1})
\qquad\text{for all } i.
\end{equation}
Applying the chain rule of probability and then \eqref{eq:markov-property} repeatedly,
\begin{empheq}[box=\fbox]{equation}
\label{eq:chain-factorisation}
\Pz(a) \;=\; \mu(a_1) \prod_{i=2}^{\nsites} \kernel(a_i \mid a_{i-1}),
\end{empheq}
where $\mu$ is the law of the first site and $\kernel$ is the \emph{transition kernel}. We
take $\kernel$ homogeneous --- the same rule at every step --- which will matter in
Chapter~\ref{ch:bp} and again when we estimate it.

\begin{intuition}
The factorisation is the whole point. A general density on $\Reals^{\nsites}$ is an
unmanageable object. Equation~\eqref{eq:chain-factorisation} says this one is built from
$\nsites-1$ copies of a single two-variable function. Everything tractable about the model
descends from that.
\end{intuition}

\section{The specific family used here}

We use additive transitions with a linear autoregression,
\begin{equation}
\label{eq:ar1}
a_i \;=\; \corr\, a_{i-1} + \varepsilon_i,
\qquad
\varepsilon_i \sim \innov \text{ i.i.d.},
\qquad
\kernel(a'\mid a) = \innov(a' - \corr a),
\end{equation}
with $|\corr|<1$, and $\innov$ an \emph{innovation density} of mean zero and variance $\ivar$.

\subsection{Stationarity and the covariance}

Suppose $a_1$ is drawn so that the chain is stationary, i.e.\ $\Var(a_i)=\sigma^2$ for all
$i$. Taking variances in \eqref{eq:ar1} and using independence of $\varepsilon_i$ from
$a_{i-1}$,
\[
\sigma^2 = \corr^2\sigma^2 + \ivar
\qquad\Longrightarrow\qquad
\sigma^2 = \frac{\ivar}{1-\corr^2}.
\]
Choosing $\ivar = 1-\corr^2$ therefore gives unit marginal variance. For the covariance,
multiply \eqref{eq:ar1} by $a_{i-k}$ and take expectations; since $\varepsilon_i$ is
independent of everything earlier,
\[
\Cov(a_i, a_{i-k}) = \corr\,\Cov(a_{i-1},a_{i-k}) = \dots = \corr^{k}\sigma^2 ,
\]
so with the normalisation above
\begin{empheq}[box=\fbox]{equation}
\label{eq:ar1-cov}
(\Sigma_0)_{ij} \;=\; \Cov(a_i,a_j) \;=\; \corr^{\,|i-j|}.
\end{empheq}

\begin{intuition}
Correlation decays geometrically with separation, with $\corr$ setting the range. At
$\corr=0.85$, sites five apart still share $\corr^5\approx0.44$ of their variation; sites
twenty apart share $0.04$. There is a well-defined correlation length, and it will reappear in
Chapter~\ref{ch:gaussian} as the scale controlling how much context a denoiser needs.
\end{intuition}

\subsection{The crucial degree of freedom}
\label{sec:matched-covariance}

Here is the feature that makes this family a good experimental probe, and it is worth
dwelling on.

Equation~\eqref{eq:ar1-cov} depends on $\innov$ \emph{only through its variance} $\ivar$.
Everything else about the innovation --- whether it is Gaussian, heavy-tailed, bounded,
bimodal --- leaves no trace in the covariance whatsoever.

So one can build a family of data distributions that are
\begin{itemize}[nosep]
\item \emph{identical} in every second-order statistic, and
\item \emph{different} in every higher-order one.
\end{itemize}

\begin{center}
\small
\begin{tabular}{@{}llcc@{}}
\toprule
family & innovation & excess kurtosis & covariance \\
\midrule
\texttt{GaussianAR1} & $\normal{0}{\ivar}$ & $0$ & $\corr^{|i-j|}$ \\
\texttt{LaplaceAR1} & Laplace$(0,b)$, $2b^2=\ivar$ & $+3$ & $\corr^{|i-j|}$ \\
\texttt{StudentTAR1} & scaled $t_\nu$ & $6/(\nu-4)$ & $\corr^{|i-j|}$ \\
\texttt{UniformAR1} & uniform on $[-c,c]$ & $-1.2$ & $\corr^{|i-j|}$ \\
\texttt{GaussianMixtureAR1} & symmetric two-component & $<0$, bimodal & $\corr^{|i-j|}$ \\
\bottomrule
\end{tabular}
\end{center}

\begin{intuition}
This is the experimental design in one table. Any method that uses only second-order
information --- and Chapter~\ref{ch:gaussian} shows that a very standard approximation does
exactly that --- cannot tell these five apart, by construction. Any difference in performance
across the row is therefore attributable to higher-order structure and nothing else. There is
no confound to argue about.
\end{intuition}

\begin{codebox}
\coderef{src/priors.py}{GaussianAR1}, \texttt{LaplaceAR1}, \texttt{StudentTAR1},
\texttt{UniformAR1}, \texttt{GaussianMixtureAR1}. Each is a frozen dataclass holding only
$\corr$, with \texttt{q}, \texttt{sample}, \texttt{log\_transition\_matrix} and
\texttt{innovation\_excess\_kurtosis} as properties. The variance matching of
\S\ref{sec:matched-covariance} is built into each class rather than imposed by the caller,
so the families cannot accidentally be compared at mismatched covariance.
\end{codebox}

\section{Why this is a reasonable model of anything}

Two defences, and it is worth being honest that they are defences rather than derivations.

\paragraph{Sequential data has temporal coherence.} Diffusion models are increasingly applied
to video and other sequential data, where consecutive frames are strongly dependent. A
first-order Markov chain is the simplest non-trivial model of that dependence. It is not a
model \emph{of} video; it is a model of the property that makes video different from a bag of
independent images.

\paragraph{It is a controllable probe.} The stronger justification is methodological. The set
of distributions for which the population score is computable is small, and every member of it
is a modelling compromise. This one buys genuine nonlinearity of the score --- which
Gaussians do not have --- at the price of a strong structural assumption. Given that the aim
is to study how estimators exploit structure, having structure one can dial is the point.

\begin{pitfall}
The model assumes \emph{exact} first-order Markov structure. If the data has longer-range
dependence, an estimator built on \eqref{eq:chain-factorisation} is misspecified, and no
amount of data repairs a misspecification. This is the single largest limitation of everything
that follows, and it is not mitigated anywhere in this document.
\end{pitfall}

\section{Sampling from the chain}

Generating data is direct: draw $a_1$ from $\mu$, then iterate \eqref{eq:ar1}. There is no
rejection, no burn-in, no approximation.

\begin{codebox}
Each prior's \texttt{sample(rng, n)} returns \emph{one} sequence of length \texttt{n}, built
by the recursion above with $a_1$ drawn standard normal. Batches are formed by stacking, e.g.
\coderef{experiments/exp\_16\_sampling\_validation.py}{sample\_chains}. The single-sequence
signature is a deliberate simplicity, not an oversight; batching is the caller's business.
\end{codebox}

\begin{remark}
Setting $a_1\sim\normal{0}{1}$ makes the chain exactly stationary only for the Gaussian family,
where the stationary law is $\normal{0}{1}$. For the others the first site is very slightly off
the stationary distribution and the chain relaxes to it within a few sites. Since $\corr=0.85$
gives a relaxation of $\corr^k$, the effect is below $10^{-3}$ by site $\sim40$, and all
reported statistics either use interior sites or average over the sequence. The alternative ---
sampling $a_1$ from the exact stationary law of each family --- would require a fixed-point
computation per family for no measurable gain.
\end{remark}
